\documentclass[10pt,a4paper]{amsart}
\usepackage[margin=19mm]{geometry}
\usepackage{amsmath,amssymb,mathtools}
\usepackage[scr=rsfs,cal=euler]{mathalfa}
\usepackage{xfrac}
\usepackage[unicode,hidelinks,bookmarks=false]{hyperref}
\newcommand{\ie}{i.e.\ }
\newcommand{\cf}{cf.\ }
\newcommand{\resp}{respectively\ }
\newcommand{\Subsection}{\subsection}
\providecommand{\nobreakdash}{\nobreak\hbox{-}}
\setlength{\emergencystretch}{2em}
\multlinegap0pt
\usepackage{amssymb}
\usepackage{mathtools}
\usepackage{bm}
\usepackage[shortlabels]{enumitem}
\usepackage{float}
\usepackage{algorithm}\usepackage{algpseudocode}
\usepackage{tikz}
\newtheorem{assumption}{Assumption}
\newtheorem{lemma}{Lemma}
\title{Homogeneous Second-Order Descent Method: Generalization and Application}
\author{Yuntian Jiang, Chuwen Zhang, Bo Jiang, Yinyu Ye}
\date{Source: arXiv:2511.00680v3 (CC BY 4.0)}
\begin{document}
\maketitle

\noindent\emph{Source and licence.} Homogeneous Second-Order Descent Method: Generalization and Application. Version arXiv:2511.00680v3 (CC BY 4.0), distributed by arXiv under the Creative Commons Attribution 4.0 International licence, http://creativecommons.org/licenses/by/4.0/, as recorded in the arXiv licence metadata for this exact version and shown on the abstract page https://arxiv.org/abs/2511.00680v3. Full licence notice: \texttt{licenses/arxiv-2511.00680-CC-BY-4.0.txt}.\par\smallskip

\noindent\textit{Editorial scope.} Counted unit: the article's lemma bound (source0.tex lines 7972--8030). The article's complete original proof of it is delivered with it (source0.tex lines 8031--8220, one \texttt{proof} environment of 190 lines). No mathematics is written, repaired or re-derived: the statement and the proof are the article's own lines, in the article's own notation, and no retained line is substituted or re-flowed.  Retained uncounted as the source-local context the proof needs: lines 151--153; lines 251--257; lines 258--264; lines 1096--1099; lines 7865--7867; lines 7925--7940.  Nothing was omitted from any retained span, so the package carries no omission record.  No retained line contains a citation, so the wrapper carries no bibliography and no bibliography entry.  No retained line contains a figure environment, an image-inclusion command or drawing code, so no image is delivered and the package declares no asset dependency.  Editorial formatting only, in addition: an AMS \texttt{amsart} preamble that restates the article's own declarations and macros used by the retained lines, namely the environments \texttt{assumption}, \texttt{lemma} on separate counters, together with the standard packages those lines need.  The retained environments print in this document as Assumption 1, Lemma 1 in order of appearance, whereas the article prints the same items with the numbering of its own full document.  The complete native source of the article as distributed in the arXiv e-print for this exact version is bundled unchanged as \texttt{sources/188270/source0.tex}.
\begin{equation} \label{eq.main problem}
    f^* = \min_{x\in \mathbb{R}^n} f(x),
\end{equation}

\begin{assumption}
    \label{assm.lipschitz}
    The objective function $f: \mathbb{R}^n \to \mathbb{R}$ is twice continuously differentiable, and its Hessian is Lipschitz continuous. That is, there exists a constant $M > 0$ such that for all $x, y \in \mathbb{R}^n$,
    \begin{equation}
        \label{eq.main assm}
        \|\nabla^2 f(x) - \nabla^2 f(y)\| \leq M \|x - y\|. \end{equation}
\end{assumption}

\begin{assumption}
    \label{assm.solvable}
    The problem \eqref{eq.main problem} is solvable; that is, there exists $x^* \in \mathbb{R}^n$ such that
    \begin{equation*}
        f(x^*) = f^*:= \min_{x \in \mathbb{R}^n} f(x).
    \end{equation*}
\end{assumption}

    \begin{equation}
        \label{eq.bracket points}
        \sigma_- = \sqrt{\frac{2M\epsilon_g}{1+2\theta}}, \quad \sigma_+ = \sqrt{\frac{MG_0}{\eta}}.
    \end{equation}

    \begin{equation}\label{eq.bounded assm point}
        \|x-x^*\|\leq M_0, \ \|v-x^*\|\leq M_0,
    \end{equation}

    \begin{equation}
        \label{eq.bounded gap bisection sigma}
        \psi_- - \psi_+
        \leq
        \frac{\sigma_+}{\sigma_-^2}
        \Bigg[
            \frac{2M_0}{\sigma_-}
            +
            \frac{4MM_0^2}{\sigma_-^2}
            +
            2\left(
            \frac{2MM_0}{\sigma_-}+1
            \right)\psi_+
            \Bigg]
        (\sigma_+-\sigma_-).
    \end{equation}

\begin{lemma}
    \label{coro.bisection bound}
    Suppose Assumption~\ref{assm.lipschitz} and
    Assumption~\ref{assm.solvable} hold, and there exists \(M_0>0\) such that
    \eqref{eq.bounded assm point} holds. Then in \textnormal{R\&B}, we have
    \begin{equation}
        \label{eq.bisection-coefficient-bound}
        \frac{\sigma_+}{\sigma_-^2}
        \left[
            \frac{2M_0}{\sigma_-}
            +
            \frac{4MM_0^2}{\sigma_-^2}
            +
            2\left(
            \frac{2MM_0}{\sigma_-}+1
            \right)\psi_+
            \right] \leq
        C(G_0,M_0,M,\epsilon_g),
    \end{equation}
    where
    \begin{equation}
        \label{eq.definition C bisection}
        C(G_0,M_0,M,\epsilon_g)
        :=
        \sqrt{G_0}\Bigg(
        \frac{(1+2\theta)^3M_0}
            {\sqrt{2}M\eta^2\epsilon_g^{3/2}}
        +
        \frac{(1+2\theta)^2M_0^2}
            {\sqrt{M}\eta\epsilon_g^2}
        +
        \frac{\sqrt{2}(1+2\theta)^{3/2}M_0}
            {M\epsilon_g^{3/2}}
        +
        \frac{(1+2\theta)^2}
            {M^{3/2}\eta^{3/2}\epsilon_g}
        \Bigg).
    \end{equation}
    If the bisection procedure does not terminate, then
    \begin{equation}
        \label{eq.bound value by brackets}
        \sigma_+-\sigma_-
        >
        \frac{1-\eta}
        {M C(G_0,M_0,M,\epsilon_g)}.
    \end{equation}
    Further, the number of oracle calls during the bisection is bounded by
    \begin{equation}
        \label{eq.bisection-oracle-bound}
        \log\left(
        \frac{
                M\sqrt{\frac{MG_0}{\eta}}\,
                C(G_0,M_0,M,\epsilon_g)
            }{
                1-\eta
            }
        \right).
    \end{equation}
\end{lemma}

\begin{proof}
    Let
    \(
    \underline{\sigma}
    :=
    \sqrt{\frac{2M\epsilon_g}{1+2\theta}},
    \overline{\sigma}
    :=
    \sqrt{\frac{MG_0}{\eta}}.
    \)
    By the construction of the bracket points in \eqref{eq.bracket points},
    throughout the bisection we have
    \(
    \underline{\sigma}
    \le
    \sigma_-
    \le
    \sigma_+
    \le
    \overline{\sigma}.
    \)
    Moreover, the upper bracket satisfies
    \(
    \psi_+\le \frac{\eta}{M}.
    \)
    Using \eqref{eq.bounded gap bisection sigma}, we get
    \[
        \psi_- - \psi_+
        \le
        B_k(\sigma_+-\sigma_-),
    \]
    where
    \[
        B_k
        :=
        \frac{\sigma_+}{\sigma_-^2}
        \left[
            \frac{2M_0}{\sigma_-}
            +
            \frac{4MM_0^2}{\sigma_-^2}
            +
            2\left(
            \frac{2MM_0}{\sigma_-}+1
            \right)\psi_+
            \right].
    \]
    We now upper bound \(B_k\). First,
    \[
        \frac{2\sigma_+M_0}{\sigma_-^3}
        \le
        \frac{
            2\sqrt{\frac{MG_0}{\eta}}M_0
        }{
            \left(
            \sqrt{\frac{2M\epsilon_g}{1+2\theta}}
            \right)^3
        }  =
        \frac{
            \sqrt{G_0}(1+2\theta)^{3/2}M_0
        }{
            \sqrt{2}M\sqrt{\eta}\epsilon_g^{3/2}
        } \le
        \frac{
            \sqrt{G_0}(1+2\theta)^3M_0
        }{
            \sqrt{2}M\eta^2\epsilon_g^{3/2}
        }.
    \]
    Second,
    \[
        \frac{4\sigma_+MM_0^2}{\sigma_-^4}
        \le
        \frac{
            4\sqrt{\frac{MG_0}{\eta}}MM_0^2
        }{
            \left(
            \sqrt{\frac{2M\epsilon_g}{1+2\theta}}
            \right)^4
        } =
        \frac{
            \sqrt{G_0}(1+2\theta)^2M_0^2
        }{
            \sqrt{M}\sqrt{\eta}\epsilon_g^2
        } \le
        \frac{
            \sqrt{G_0}(1+2\theta)^2M_0^2
        }{
            \sqrt{M}\eta\epsilon_g^2
        }.
    \]
    Third, using \(\psi_+\le \eta/M\),
    \[
        \frac{4\sigma_+MM_0\psi_+}{\sigma_-^3}
        \le
        \frac{4\eta M_0\sigma_+}{\sigma_-^3} \le
        \frac{
            \sqrt{2G_0}(1+2\theta)^{3/2}M_0\sqrt{\eta}
        }{
            M\epsilon_g^{3/2}
        } \le
        \frac{
            \sqrt{2G_0}(1+2\theta)^{3/2}M_0
        }{
            M\epsilon_g^{3/2}
        }.
    \]
    Finally,
    \[
        \frac{2\sigma_+\psi_+}{\sigma_-^2}
        \le
        \frac{2\eta\sigma_+}{M\sigma_-^2} \le
        \frac{
            \sqrt{G_0}(1+2\theta)\sqrt{\eta}
        }{
            M^{3/2}\epsilon_g
        } \le
        \frac{
            \sqrt{G_0}(1+2\theta)^2
        }{
            M^{3/2}\eta^{3/2}\epsilon_g
        }.
    \]
    Combining the above four estimates gives
    \[
        B_k
        \le
        C(G_0,M_0,M,\epsilon_g),
    \]
    where \(C(G_0,M_0,M,\epsilon_g)\) is defined in
    \eqref{eq.definition C bisection}. Hence
    \[
        \psi_- - \psi_+
        \le
        C(G_0,M_0,M,\epsilon_g)(\sigma_+-\sigma_-).
    \]

    If the bisection procedure does not terminate, then by the bracket condition
    we have
    \(
    \psi_- - \psi_+
    >
    \frac{1-\eta}{M}.
    \)
    Therefore,
    \[
        \sigma_+-\sigma_-
        >
        \frac{1-\eta}
        {M C(G_0,M_0,M,\epsilon_g)}.
    \]
    This proves \eqref{eq.bound value by brackets}.

    By the mechanism of bisection, the total number of bisection steps in the
    \(k\)-th iteration, denoted by \(N_k\), satisfies
    \[
        N_k
        \le
        \log\left(
        \frac{
                \left(
                \sqrt{\frac{MG_0}{\eta}}
                -
                \sqrt{\frac{2M\epsilon_g}{1+2\theta}}
                \right)
                M C(G_0,M_0,M,\epsilon_g)
            }{
                1-\eta
            }
        \right)
        \le
        \log\left(
        \frac{
                M\sqrt{\frac{MG_0}{\eta}}\,
                C(G_0,M_0,M,\epsilon_g)
            }{
                1-\eta
            }
        \right).
    \]
    Since \(G_0\) and \(M_0\) are polynomially bounded in \(D_0\), by the
    definition of \(C(G_0,M_0,M,\epsilon_g)\), and omitting fixed algorithmic
    parameters, we conclude
    \(
    N_k
    =
    O\left(
    \log\frac{MD_0}{\epsilon_g}
    \right).
    \)
\end{proof}

\end{document}
